\documentclass[10pt]{article}
\usepackage[margin=0.75in]{geometry}
\usepackage{amsmath,amssymb,amsthm,microtype}
\newtheorem{theorem}{Theorem}
\title{Noncontraction Transport Through Unequal-Length Merges}
\author{Collatz proof project}
\date{September 5, 2026}
\begin{document}
\begin{center}
{\Large Noncontraction Transport Through Unequal-Length Merges}\par\medskip
Collatz proof project\quad---\quad September 5, 2026
\end{center}
\begin{abstract}
An unequal-length merging predecessor inherits full noncontraction when
its finite prefix is noncontracting and its final coefficient dominates
the original prefix coefficient. We formalize this transport principle
and an infinite arithmetic family meeting after one versus twelve steps.
A bounded exact search supplies additional examples but leaves a global
existence obligation. Neither nondivergence nor Collatz is proved.
\end{abstract}
Let $T(n)=n/2$ for even $n$ and $(3n+1)/2$ for odd $n$. Write
$j_t(n)$ for the odd count and $C_t(n)=3^{j_t(n)}/2^t$.
Call $n$ noncontracting if $C_t(n)\geq1$ for every $t\geq0$.

\begin{theorem}[Dominating merge]
Suppose $n$ is noncontracting, $T^s(x)=T^t(n)$,
$C_i(x)\geq1$ for $0\leq i\leq s$, and $C_s(x)\geq C_t(n)$.
Then $x$ is noncontracting.
\end{theorem}
\begin{proof}
For $v\geq0$, the coefficient cocycle and the shared endpoint give
\[
C_{s+v}(x)=C_s(x)C_v(T^s(x))
 =\frac{C_s(x)}{C_t(n)}C_{t+v}(n)\geq1.
\]
Earlier times are covered by the prefix hypothesis. The Lean proof uses
natural-number cross multiplication and cancellation of a positive power
of three. It does not assume the shifted seed $T^t(n)$ is noncontracting;
that property does not follow just by shifting the original hypothesis.
\end{proof}

\begin{theorem}[An infinite merging family]
For every integer $Q\geq0$, put
\[
N=111+4374Q,\qquad M=103+4096Q.
\]
Then $0<M<N$, and
\[
T(N)=T^{12}(M)=167+6561Q,\qquad j_1(N)=1,\qquad j_{12}(M)=8.
\]
Every prefix of the twelve-step path from $M$ is noncontracting.
Consequently, if $N$ is noncontracting, so is the smaller seed $M$.
\end{theorem}
\begin{proof}
Lean's kernel verifies $T^{12}(103)=167$, $j_{12}(103)=8$ and all
thirteen prefix coefficient inequalities. Binary-cylinder transport
preserves these parities for $103+2^{12}Q$ and adds $3^8Q$ to its
endpoint. The seed $N$ is odd, giving the same endpoint after one step.
The coefficient comparison is
\[
\frac{C_{12}(M)}{C_1(N)}
=\frac{3^8/2^{12}}{3/2}=\frac{2187}{2048}>1.
\]
The first theorem applies, and $N-M=8+278Q>0$.
\end{proof}
Thus no least positive noncontracting seed can belong to this progression.
All $N$ are multiples of three, and the meeting occurs after an increasing
forward step. The result is not a construction of a noncontracting seed,
nor an exclusion of noncontraction for every member without the minimality
condition. No mathematical priority claim is made.

\paragraph{Finite search.}
The exact Python search takes the 7,495 seeds $0<n<2^{18}$ whose first
18 prefix coefficients are at least one. Equal-time/equal-count replacement
already covers 1,391 of them as ordinary smaller merges. Searching both
forward and replacement lengths from zero through 18 covers another 3,768.
All 3,768 additions have certificates satisfying the first theorem;
937 of the selected certificates use a positive forward time. Each is
independently replayed, and no original path descends before its meeting.
The 2,336 unresolved seeds begin $27,327,447,495,639,667,703,763$.
These are finite seed counts, not coverage of all their cylinder lifts.
Failure within these bounds says nothing about longer replacements.

\paragraph{Verification and remaining gap.}
The three public declarations in \texttt{Collatz.DominatingMerge} are
included in the axiom audit. The reproducible search is
\texttt{analysis/variable\_length\_surgery.py}; its tests compare against
every smaller seed at depth eight and replay five members of the family.
The census is not kernel certified. A proof must still force a suitable
smaller predecessor for every hypothetical least noncontracting seed;
the present certificates cover only particular cases. Nontrivial cycles
also remain outside the conclusion.

\newpage
\section*{Kernel-certified arithmetic progressions}
The individual merging examples can be turned into rules that apply to
arbitrarily large seeds. This extension certifies their validity, not
universal applicability.

\begin{theorem}[Progression lifting]
Let $n,x,t,s,a,b$ be natural numbers satisfying $0<x<n$,
$T^s(x)=T^t(n)$, $C_i(x)\geq1$ for $0\leq i\leq s$, and
$C_s(x)\geq C_t(n)$. Write $j=j_t(n)$ and $k=j_s(x)$, and assume
\[
3^ja=3^kb,\qquad 2^sb\leq2^ta.
\]
For every $Q\geq0$, put $N_Q=n+2^taQ$ and $X_Q=x+2^sbQ$.
Then $0<X_Q<N_Q$. If $N_Q$ is noncontracting, so is $X_Q$.
\end{theorem}
\begin{proof}
Binary-cylinder transport gives
\[
T^t(N_Q)=T^t(n)+3^jaQ,
\qquad T^s(X_Q)=T^s(x)+3^kbQ.
\]
The balance condition makes these endpoints equal. Prefix parities and
odd counts through $t$ and $s$ are unchanged, so the noncontracting-prefix
and coefficient-domination checks remain valid. The first theorem
therefore transfers full noncontraction. Finally, $x<n$ and $2^sb\leq2^ta$
give strict ordering for every nonnegative quotient, while $x>0$ gives
positivity.
\end{proof}

\paragraph{Choosing the steps.}
A reduced choice is
\[
a=3^{\max(k-j,0)},\qquad b=3^{\max(j-k,0)}.
\]
It satisfies the balance equation. Coefficient domination gives the
required ordering of the two seed increments. The generator shifts a
certificate backward along this progression only while both seeds stay
positive and the replacement remains strictly smaller. The shifted
certificate is checked again, including every replacement-prefix parity.
These normalization operations are not assumed correct by Lean: the
resulting finite data must independently satisfy its validator.

\paragraph{What the kernel checks.}
\texttt{MergeRule.Valid} records positive seed ordering, the exact common
endpoint, every replacement-prefix coefficient inequality, coefficient
domination, ternary balance, and ordered seed increments.
\texttt{MergeRule.sound} proves the theorem above for any valid rule.
The generated \texttt{MergeRuleTable.rules} contains 360 rules. Its
\texttt{rules\_valid} theorem checks them with ordinary kernel reduction
using \texttt{decide}; no external verifier result or added axiom supplies
validity. \texttt{certified\_rule\_sound} applies the generic theorem to
any listed rule and any natural quotient.

\paragraph{Scope of the finite generation.}
The 3,768 successful dominating certificates from the preceding search
normalize to 360 distinct progressions. The inclusion filter removes
no further distinct progression in this run. The listed rules include
the earlier $111+4374Q$ family and simple inverse families such as
$2+3Q$ merging with $1+2Q$. The finite Python search and its completeness
remain separate from the kernel-certified validity of the resulting list.

Every listed progression is excluded for a least positive noncontracting
seed: its certified replacement would be smaller and noncontracting.
No theorem says that this finite union contains every such hypothetical
seed. In particular, the earlier unresolved seeds do not acquire an
in-budget certificate from this normalization. The missing existence
argument may need a seed-dependent, unbounded family of rules. The
nontrivial-cycle problem is also unchanged.

\paragraph{Reproduction.}
\begingroup\raggedright
Run \texttt{python3 analysis/merge\_progressions.py --write-lean} from the
root, then \texttt{lake build Collatz.MergeRuleTable} from \texttt{lean/}.
The generator's JSON and the Lean source are both stored in the repository.
Tests check normalization, progression inclusion with lower bounds, and
exact agreement between the generated source and the saved rule data.
The theorem audit covers the generic soundness theorem and both public
table theorems. No mathematical priority claim is made.\par\endgroup
\end{document}
